\section*{Conclusion}

We examined whether making a multi-UAV delivery policy tail-risk sensitive makes the operation more resilient to distribution shift. To answer the question under controlled conditions we built a stochastic urban delivery environment in which wind, navigation quality, traffic and energy can each be shifted independently, trained six policies with three seeds each under a matched budget, and evaluated them on paired scenarios with every policy flown under every return-to-home reserve rule.

Two conclusions follow. First, tail-risk objectives did not deliver a reliable reduction in catastrophic losses. A fixed conditional value-at-risk objective, the same objective with a risk level scheduled online from collision threat, energy margin, traffic density and disturbance intensity, and a variant that places the tail measure on a learned catastrophe-cost distribution all finished within the bootstrap and seed-to-seed spread of a risk-neutral soft actor-critic, and a more cautious deployment risk level did not lower collision rates. The data point to three reasons: the lower tail of the return is dominated by ordinary variability in time and energy rather than by collisions, caution expresses itself as hesitancy that lengthens exposure to conflicts, and a large share of the residual risk comes from aircraft the policy never observes. Second, statistical calibration of the return-to-home decision was decisive for energy safety. The online conformal reserve held its specified coverage of return-leg energy in every shift family, \RCrestcOnlineCovEnergy\% under the energy shift where the conventional fixed margin reached only \RCrestcStdCovEnergy\%, and it cut energy depletion under that shift from \RCrestcStdDepEnergy\% to \RCrestcOnlineDepEnergy\% of missions. The price is paid in abandoned deliveries, and the frontier we report lets an operator choose the exchange rate explicitly instead of inheriting it from a fixed margin.

The broader lesson is to place statistical guarantees where the failure mechanism is modelled and a conservative fallback always exists, and to treat claims of resilience from tail-risk objectives as hypotheses to be tested against strong risk-neutral baselines on paired, family-by-family shifts. The same reserve construction, a modelled quantity with a monotone knob and an always-available backup action, transfers naturally to separation assurance with calibrated bounds on miss distance, which we see as the most promising route to reducing the conflict losses that neither mechanism studied here addressed.
